\section{Dominio de dependencia numérico. Condición CFL}

Recordamos que la ecuación \eqref{eq:adveccion 1d} tiene por solución $u(x,t) = u_0(x - at)$. Esto quiere decir que para conocer $u$ en el punto $(x,t)$, me basta con conocer el valor de $u$ en el punto $(x-at, 0)$. O en cualquier punto de la forma $(x-a(t-s),s)$.
Llamamos \emph{dominio de dependencia} a esta región espacial.
\begin{example}[Dominio de dependencia de \eqref{eq:adveccion 1d}]
    \begin{equation*}
        \Big\{ (x - a(t-s) , s) : s \in (0,t) \Big\}.
    \end{equation*}
\end{example}

\begin{example}[Dominio de dependencia de \eqref{eq:wave} en $d=1$]
    La solución explícita es \eqref{eq:wave 1d D'Alambert}, de donde deducimos que
    el dominio de dependencia del punto $(x,t)$ es el triángulo
    $$
        \Big\{ (y,s) : s \in (0,t), x-c (t-s) \le y \le x + c(t-s) \Big\} .
    $$
\end{example}
\todo{Elegir definiciones rigurosas}
En el caso de \eqref{eq:adveccion 1d} y \eqref{eq:adveccion 1d EE backward}, vemos que \eqref{eq:adveccion 1d CFL} dice precisamente esto. Además, hemos comprobado que es condición suficiente para la estabilidad del método.
\todo{figura}

\begin{remark}[Condición Courant-Friedrichs-Lewy o CFL]
    Un método numérico sólo puede ser convergente si el dominio de dependencia numérico contiene al dominio de dependencia la EDP, al menos cuando $\tau$ y $h$ tienden a $0$.
\end{remark}




